\section{Conclusion and perspectives}
\label{sec:conclusion}



The objective of this internship was to study how IO500 can be tuned
without exploring the complete parameter space with valid benchmark
executions. A valid IO500 measurement is expensive because some phases
must run for at least 300 seconds, and a complete execution can take
around 30 minutes. As a result, testing every possible configuration with
the official settings would require a very large amount of computing
time.

To make the study possible, a reduced-cost experimental setup was used.
The stonewall time was reduced to 10 seconds and the size-related
parameters were reduced. This made it possible to explore the parameter
space much faster and to build a reference dataset containing the IO500
SCORE obtained for the tested configurations.

This reference dataset was then used to compare different search
strategies. Random Search was used as a simple baseline, while the main
part of the study focused on Bayesian Optimization. Several Bayesian
Optimization setups were compared by changing the surrogate model, the
acquisition function, the exploration parameters, and the number of
initial points.

The results show that the choice of Bayesian Optimization setup has an
important effect on performance. Gaussian Process gave the best results
among the surrogate models considered in this study. The best setups used
either LCB or EI as acquisition functions, with only small differences
between them at the top of the ranking.

A well-chosen Bayesian Optimization configuration reached approximately
98\% of the best SCORE available in the reference dataset with a limited
number of evaluations. In comparison, Random Search had more difficulty
reaching the very small region containing the highest scores.

The comparison with a human specialist also showed that expert knowledge
can quickly identify a good configuration. The best specialist
configuration reached approximately 95\% of the optimum, while Bayesian
Optimization reached approximately 98\%. This comparison should however
be interpreted carefully because it was based on only one specialist and
a small number of proposed configurations.

The internship also contributed to the development of IOPS. IO500 support
was added through a configuration generator, allowing users to create an
IO500 YAML template with one \add{simple} command. The IOPS documentation was updated to describe this functionality and the
main parameters that can be adapted to the target system. The search
component was also extended with an implementation of the
Tree-structured Parzen Estimator (TPE). TPE was implemented during the
internship but was not evaluated in the IO500 experiments.

Overall, this work shows that guided search is a promising approach for
IO500 tuning. Instead of testing a large number of configurations without
guidance, Bayesian Optimization can use previous observations to focus
the search on more promising parts of the parameter space.



Several directions can extend the work presented in this report. The first step is to complete the validation of the best configurations
using valid IO500 executions. This would make it possible to confirm the
results obtained with the reduced-cost setup and to report official-grade
IO500 scores. A second direction is to reproduce the study on IRENE at TGCC. IRENE is a
larger and more complex HPC system, so repeating the experiments there
would help determine whether the conclusions obtained on PlaFRIM can also
be observed on another platform.

The 10-node challenge could also be extended. This would provide an
additional experimental setting for studying how the search methods behave
when the available computing resources are limited. The comparison with human expertise could also be improved. In the current
study, only a small number of configurations were provided by one
specialist. Collecting configurations from more specialists would provide
a stronger human baseline and make the comparison more meaningful.

Another direction is to add and compare more search methods. Possible
approaches include genetic algorithms, simulated annealing, and
greedy or coordinate-based search. In particular, implementing a
coordinate search would make it possible to reproduce the approach used
for the CLAIX23 study directly on the dataset produced during this work.

The results of the study could also be used to improve the
\texttt{BayesianPlanner} in IOPS. For example, the default surrogate model
could be selected from the results obtained in this study. The handling of
discrete parameters and the behaviour of the planner when it proposes an
already evaluated configuration could also be improved.

Finally, a longer-term direction is to study transfer learning between HPC
systems. Instead of starting the optimization from zero on every machine,
information collected on one system could be reused to initialize the
search on another system. This could further reduce the number of IO500
executions required to find a good configuration.